\documentclass[11pt]{article}

% basic packages
\usepackage[margin=1in]{geometry}
\usepackage[pdftex]{graphicx}
\usepackage{amsmath,amssymb,amsthm}
\usepackage{custom}
\usepackage{lipsum}
\usepackage{booktabs}
\usepackage{enumitem}

% page formatting
\usepackage{fancyhdr}
\pagestyle{fancy}

\renewcommand{\sectionmark}[1]{\markright{\textsf{\arabic{section}. #1}}}
\renewcommand{\subsectionmark}[1]{}
\lhead{\textbf{\thepage} \ \ \nouppercase{\rightmark}}
\chead{}
\rhead{}
\lfoot{}
\cfoot{}
\rfoot{}
\setlength{\headheight}{14pt}

\linespread{1.03} % give a little extra room
\setlength{\parindent}{0.2in} % reduce paragraph indent a bit
\setcounter{secnumdepth}{2} % no numbered subsubsections
\setcounter{tocdepth}{2} % no subsubsections in ToC

\begin{document}

% make title page
\thispagestyle{empty}
\bigskip \
\vspace{0.1cm}

\begin{center}
{\fontsize{22}{22} \selectfont Lecture Notes on}
\vskip 16pt
{\fontsize{36}{36} \selectfont \bf \sffamily General Relativity}
\vskip 24pt
{\fontsize{18}{18} \selectfont \rmfamily Owen Lin} 
\vskip 6pt
{\fontsize{14}{14} \selectfont \ttfamily owen2006lin@gmail.com} 
\vskip 24pt
\end{center}

{\parindent0pt \baselineskip=15.5pt
These notes cover general relativity at the undergraduate level. The sources used are:
\begin{itemize}
    \item Schutz, \textit{A First Course in General Relativity}
    \item Carroll, \textit{Spacetime and Geometry}
\end{itemize}


Potential sources to consult in the future:

\begin{itemize}
    \item Hartle, \textit{Gravity: An Introduction to Einstein's General Relativity}
    \item Misner, Thorne, Wheeler, \textit{Gravitation}
    \item Weinberg, \textit{Gravitation and Cosmology}
    \item Wald, \textit{General Relativity}
\end{itemize}

}
 


% make table of contents
\newpage
\microtoc
\newpage

% main content
\section{Preliminaries}

\subsection{Coordinate Transformations}
First, we'll briefly recap special relativity using index notation. 
    \begin{itemize}
        \item Let's define a four vector \( x^{\mu} \) where \( x^{0} \) is the temporal component and \( x^{i} \) (\( i \neq 0 \)) are the spacial coordinates.
        In these notes, we'll write the representation of a vector in coordinate transformation as \( x^{\mu'} \) (each source has     their own convention). 
        
        \item  The matrix to transform the vector \( x^{\nu} \) into the new primed coordinates is
        \[ 
            x^{\mu'} = \Lambda^{\mu'}{}_\nu x^{\nu}
        \]
        An object that transforms like vector components is said to transform \textbf{contravariantly}. Note we are using the Einstein convention where summations over repeated indices drop
        the \( \Sigma \). 
    
        \item A Lorentz Transformation is one that keeps the spacetime interval constant and fix the origin so this requires \( \eta = \Lambda^{T} \eta \Lambda \). In index notation, this is:
        \[ 
            \eta_{\alpha \beta} = \Lambda^{\mu'}{}_{\alpha} \Lambda^{{\nu'}}{}_{\beta}\eta_{\mu' \nu'}
        \]
        Note that when writing the equation using matrices, the order is important because of the limitations of matrix multiplication. However, the order is irrelevant
        in index notation. These matrices are \textbf{Lorentz Transformations} and are elements of the \textbf{Lorentz Group}.

        \item A brief aside, these matrices may seem very similar to the rotation group SO(3). The Orthogonal Group is all 3x3 matrices that satisfy \( \mathbf{R}^{\mathbf{T}}\mathbf{R} = \mathbf{1} \)
        (SO(3) then are all special matrices in O(3), ie determinant 1). The difference between O(3) and the Lorentz group is that instead of the identity matrix, it is equal
        to the metric \( \eta \). The identity matrix is simply the metric in \textbf{Euclidean Space}, and \( \eta \) is similarly the metric for a \textbf{Lorentzian Space}.
        
        \item Lorentz transformations come in two forms: rotations and boosts. Below is an example of a boost in the x direction:

        \[ 
            \Lambda^{\mu'}{}_{\nu'} = 
            \begin{pmatrix}
            \cosh \phi & -\sinh \phi & 0 & 0 \\
            -\sinh \phi & \cosh \phi & 0 & 0 \\
            0 & 0 & 1 & 0 \\
            0 & 0 & 0 & 1 \\
            \end{pmatrix}
        \]

        The transformed quantity for \( t' \) and \( x' \) are then \( t' = t\cosh\phi - x\sinh\phi \) and \( x' = -t\sinh\phi + x\cosh\phi \). The velocity follows as \( v = \tanh\phi \), and we can
        then rewrite in more familiar form 
        \[ 
            t' = \gamma(t - vx) 
        \]
        \[ 
            x' = \gamma(x - vt) 
        \]
        where \( \gamma = 1/\sqrt{1 - v^{2}}  \). From here, we can proceed and obtain time dilation, length contraction, etc.

        \item Another transformation that preserves the spacetime interval is a pure translation:
        \\
        \( x^{\mu} \rightarrow x^{\mu'} =( x^{\mu} + a^{\mu})\)
        The set of all translations form an Abelian group but generally Lorentz transforms do not commute so the Lorentz Group is nonabelian. The set of all Lorentz transforms and translations form a larger
        nonabelian group: the \textbf{Poincar\'e Group}.

    \end{itemize}

\noindent Now we'll take a look at how vectors and their components transform under a change of coordinates.
    \begin{itemize}
        \item Suppose in an arbitrary vector space we have a basis of vectors \( \hat{e}_{\mu} \). Then any arbitrary vector can be written as \( A = A^{\mu}\hat{e}_{\mu} \) where \( A^{\mu} \) are the components
        of the vector. Often we will just call \( A^{\mu} \) the vector for shorthand. 

        \item An example is the tangent vector \( V^{\mu} = \frac{dx^{\mu}}{d\lambda} \). Under a lorentz transformation the components change by \( \frac{dx^{\mu'}}{dx^{\mu}} \) by the chain rule so the components
        transform as \( V^{\mu} \rightarrow V^{\mu'} = \Lambda^{\mu'}{}_\nu V^{\nu}\) but the vector itself is invariant under transformation. So we must have 
        \(V^{\mu}\hat{e}_{\mu} = V^{\nu'}\hat{e}_{\nu'} = \Lambda^{\nu'}{}_\mu V^{\mu}\hat{e}_{\nu'} \). But this must hold for any vector no matter the components so we must have 
        \( \hat{e}_{\mu} = \Lambda^{\nu'}{}_\mu \hat{e}_{\nu'} \). This means basis vectors transform opposite to their components.

        \item Note that applying a transformation and then its inverse should result in the identity. In index notation, this means \(  \Lambda^{\mu}{}_{\nu'} \Lambda^{\nu'}{}_{\rho} = \delta^{\mu}_\rho \)

        \item We can define a vector space dual to our previous vector space, which is the set of all linear maps from the original vector space to \( \mathbb{R} \). We can write a set of basis dual vectors as 
        \( \hat{\theta}^{{\nu}}(\hat{e}_{\mu}) = \delta^{\nu}_\mu \). Then every dual vector can be written in component form \( \omega = \omega_{\mu}\hat{\theta}^{\mu}. \) (we also refer to dual vectors as "1-forms").

        \item The components of dual vectors transform like basis vectors (that is, inverse to vector components) 
        \[ 
            \omega_{\mu'} = \Lambda^{\nu}{}_{\mu'}\omega_{\nu} 
        \]
        An object that transforms like this is said to transform \textbf{covariantly}. Meanwhile, dual basis vectors transform like
        vector components \( \hat{\theta}^{\rho'} = \Lambda^{\rho'}{}_{\sigma}\hat{\theta}^{\sigma} \). The key takeaway is that the location of the index tells us how the object transforms. Upper index transform like vector
        components, while lower transforms inverse to that. Accordingly, scalars are invariant under transformation.

        \item An example is the gradient, with components \( \frac{\partial }{\partial x^{\mu}} \). Under transformation, this becomes 
        \(\frac{\partial }{\partial x^{\mu'}} = \frac{\partial x^{\mu}}{\partial x^{\mu'}} \frac{\partial }{\partial x^{\mu}} \) by the chain rule, but this is equivalent to just \( \Lambda^{\mu}{}_{\mu'}\frac{\partial}{\partial x^{\mu}} \).
        Since the gradient is a dual vector, we tend to write it with lower index as \( \partial_{\mu} \).

    \end{itemize}
\noindent We'll later generalize this idea of transformation under change of coordinates to any transformation, not just Lorentz transformations.

\subsection{Equivalence Principles}
The equivalence principles are important in that they describe gravity is a property of spacetime's geometry, not the result of any additional field.
\begin{itemize}
    \item The \textbf{Weak Equivalence Principle} states that inertial mass and gravitational mass of any object are equal.
    \[
        m_i = m_a
    \]
    This also means that objects traveling through some gravitational potential \(\Phi\) have the same acceleration \(\mathbf{a} = - \nabla \Phi\). A path through spacetime where particles don't feel any external force
    is called inertial, and all particles will locally travel this same path, a fact that isn't true for other forces (eg; such paths don't exist for EM). 
    
    \item The \textbf{Einstein Equivalence Principle} states in small regions of spacetime, the laws of physics reduce to that of special relativity and its impossible to detect the existence of gravitational field with local experiments.
    For example, a person standing on the Earth and Jupiter's surface will see differences in their experiments. But if instead this person were in a sealed container travelling uniformly through the gravitational field, there would be no way
    to discern the two apart. Note this is only true if the box covers a small, local region of spacetime. In larger regions, tidal forces come into play which can be detected.

    \item It is still possible to have a theory that violates the Einstein principle but satisfies the weak equivalence principle. For example, a theory of gravity that all particles start rotating as they travel the same path in a gravitational 
    field satisfies the weak equivalence principle but violates the Einstein principle, since you could detect a gravitational field this way.

    \item The \textbf{Strong Equivalence Principle} extends the Einstein Equivalence principle to massive particles, and particles that exert gravitational force on themselves, such as planets and black holes. 
     
    \item An important corollary of the equivalence principles is that the acceleration due to gravity is not well defined. There's no such thing as a neutral object that we could measure the acceleration with respect to. It makes more sense then, to 
    define freely falling through a gravitational field as being unaccelerated.

    \item Locality is very important to keep in mind. In special relativity we define a inertial frame by assigning every point in spacetime a clock and then connecting them with rigid rods, however this is impossible in the presence of gravity,
    because we defined freely falling objects as not accelerating, but they clearly are accelerating with respet to such a reference frame. Hence, we can only define such an inertial frame locally.   

    
\end{itemize}

\noindent \textbf{Example.} The most famous example demonstrating the equivalence principle gravitational redshift. Consider two boxes each a distance \(z\) apart in free space and nonrelativistically moving with constant acceleration \(a\) in the same direction. If we emit a 
photon from one box to the other, the photon reaches the second box at time \(\Delta t = z / c\). However, the boxes will have picked up an additional speed \(\Delta v = a \Delta t = az/c\). So the photon reaching the box will have been redshifted by
\[
    \frac{\Delta\lambda}{\lambda} = \frac{az}{c^2}
\]
We can expect then, that in a uniform gravitational field with same magnitude, we should observe the same redshift, which experiments have verified. A photon travelling down a gravitational field is redshifted after reflection.


\subsection{Manifolds}
Intuitively, a manifold is a space that locally looks Euclidean, ie looks like \( \mathbb{R}^{n} \). By ``looking Euclidean," we don't necessarily need the metric to be Euclidean, rather functions and coordinates work like they should normally. In this case, we call the dimension of
the manifold \( n \).
\\
\\
\( \mathbb{R}^{n} \) is obviously a manifold because it looks like \( \mathbb{R}^{n} \) not only locally, but globally. A sphere is of dimension \( 2 \) because locally, it looks like \( \mathbb{R}^{2} \). Meanwhile, a line conected perpendicular to a plane is not a manifold, because the 
dimension is not well defined. A cone is a manifold, however there is a singularity at the tip. Manifolds can be differentiable, but the cone has a problem at the tip, so its not a differentiable manifold. There sphere is more well behaved, so we say it's differentiable.
\\
\\
While the description of a manifold does seem intuitive, we need some heavy lifting to be able to define it more rigorously.
\begin{itemize}
    \item Let \( \phi \colon \mathbb{R}^{m} \rightarrow \mathbb{R}^{n} \) be a map of sets. We say \( \phi \) is \( C^{p} \) if the \textit{p}th derivative exists and is continuous. If a map is \( C^{\infty} \) then it is \textbf{smooth}.
    
    \item We say sets \( M \) and \( N \) are \textbf{diffeomorphic} if there exists a \( C^{\infty} \) map \( \phi \colon M \rightarrow N \) with inverse \( \phi^{-1} \colon N \rightarrow M \) that is also \( C^{\infty} \).
    
    \item An \textbf{open ball} is the set of all points \( x \) in \( \mathbb{R}^{n} \) such that \( |x - y| < r \) for some \( y \in \mathbb{R}^{n} \) and \( r \in \mathbb{R} \), where \( |x - y| \) is defined by the euclidean distance \( [\sum_{i}(x^{i} - y^{i})^{2}]^{1 / 2} \). Any open set
    in \( \mathbb{R}^{n} \) can be constructed by the union of open balls. 
    
    \item A \textbf{chart} is a subset \( U \) of some set \( M \) equipped with a one to one map \( \phi \colon U \rightarrow \mathbb{R}^{n} \) such that the image of \( \phi(U) \) is open in \( \mathbb{R}^{n} \). We can then say \( U \) is an open set in \( M \). A chart is like a coordinate system,
    we're able to describe any point in \( \mathbb{R}_{n} \) using elements in some abstract set.
    
    \item An \textbf{atlas} is a collection of charts \( \{(U_{\alpha}, \phi_{\alpha})\}  \) that satisfies two conditions. First, the union of \( U_{\alpha} \) is \( M \). Second, if any two charts overlap, \( U_{\alpha} \cap U_{\beta} \neq 0\), then the map \( \phi_{\alpha} \phi_{\beta}^{-1} \)
    maps \( \phi_{beta}(U_{\alpha} \cap U_{\beta}) \subset \mathbb{R}^{n} \) onto an open set \( \phi(U_{\alpha} \cap U_{\beta}) \subset \mathbb{R}^{n} \), where all of these maps are \( C^{\infty} \). This second condition is a way to say the charts are smoothly ``sewn together.'' If there exists
    two different coordinate systems for the same points, then transitioning between them is smooth and well defined. This allows us to describe smoothness for functions, take derivatives etc, and their behavior shouldn't depend on the coordinate system. 

    \item A \textbf{manifold} is a set \( M \) with a maximal atlas; meaning it contains every possible chart. This requirement is important, as two different atlas corresopnding to the same space shouldn't be different manifolds.
\end{itemize}

\noindent \textbf{Example.} Consider \( S^{1} \). The conventional coordinate system using angles \( \theta \colon S^{1} \rightarrow \mathbb{R} \) runs into a problem. If we include either \( \theta = 0 \) or \( \theta = 2\pi\) then we have a closed interval. But if we exclude both,
we haven't covered the whole circle. So we need at least two charts to cover the circle.
\\

\subsection{Tensor Analysis}
\noindent Now we'll describe tensors on manifolds. Unlike in flat space, vectors can't be freely moved around, so it makes more sense to describe them as associated with a single point. On a manifold, the set of all vectors associated with a single point is the tangent space to that point. 
\begin{itemize}
    \item The tangent space to a manifold is the space of directional derivative operators along curves that pass through that point. Namely, if \(f \colon M \rightarrow \mathbb{R}\) then its directional derivative is an element of the tangent space. This way, we can define tangent space without the notion of
    tangent vectors, which would be circular.

    \item A good basis to choose for the tangent space is the set of all partials \(\{\partial_\mu\}\). This is because 
    \[
        \frac{d}{d \lambda} = \frac{dx^\mu}{d\lambda}\partial_\mu   
    \]
    so the directional derivative can be expressed as a linear combination of the partials. We call \(\frac{dx^\mu}{d\lambda}\) the components of the tangent vector. This basis \(\hat{e}_\mu = \partial_\mu\) is the \textbf{coordinate basis}.

    \item Using this definition, we can see that basis vectors transform covariantly
    \[
        \partial_{\mu'} = \frac{\partial x^\mu}{\partial x^{\mu'}}\partial_\mu
    \]
    and similarly, components transform contravariantly
    \[
        V^{\mu'} = \frac{\partial x^{\mu'}}{\partial x^\mu}V^\mu
    \]

    This is more general than our definition before, because we drop the need for any Lorentz transform. Here we can define tensors under any change of basis, not just linear transformations.
    
    \item On a manifold, all vectors are tangent vectors, which means they are associated with a directional derivative. A vector field is then defines a map from smooth functions to smooth functions over the manifold, by defining a directional derivative at each point.
    
    \item Given vector fields \(X\) and \(Y\), we define the commutator \([X,Y]\) by how it acts on a function \(f(x^\mu)\)
    \[
        [X,Y]f = X(Yf) - Y(Xf)
    \]

    \item The gradient of a function is a one-form. The set of all gradients of the coordinate functions \(x^\mu\) form a basis for the cotangent space, so the basis for one-forms is \(\{dx^\mu\}\). They transform contravariantly
    \[
        dx^{\mu'} = \frac{\partial x^{\mu'}}{\partial x^\mu}dx^\mu
    \]
    and components of one forms transform covariantly
    \[
        \omega_{\mu'} = \frac{\partial x^\mu}{\partial x^{\mu'}}\omega_\mu
    \]
    

    
\end{itemize}

\noindent Finally, we can generalize our transformation laws by defining tensors. A tensor is a multilinear map of dual vectors and vectors to \( \mathbb{R} \). For example, \( T(a\omega, cV) = acT(\omega, V) \).
A tensor that takes in \( k \) vectors and \( l \) dual vectors is a type \( (k, l) \) tensor. It's easy to see that scalars are \( (0,0) \) tensors
and vectors and dual vectors are \( (1, 0) \) and \( (0, 1) \) tensors accordingly. Some important facts about tensors:

\begin{itemize}
    \item Suppose \( T \) is a \( (k,l) \) tensor, then under any general coordinate transformation \(T\) behaves as follows:
    \[
        T^{\mu_1' \cdots \mu_k'}{}_{\nu_1' \cdots \nu_l'} = \frac{\partial x^{\mu_1'}}{\partial x^{\mu_1}} \cdots \frac{\partial x^{\mu_k'}}{\partial x^{\mu_k}} \frac{\partial x^{\nu_1}}{\partial x^{\nu_1'}} \cdots \frac{\partial x^{\nu_l}}{\partial x^{\nu_l'}} T^{\mu_1 \cdots \mu_k}{}_{\nu_1 \cdots \nu_l}
    \]
    We say that \(T\) transforms \(k\) times contravariantly and \(l\) times covariantly

    \item Suppose \( T \) is a \( (k,l) \) tensor and \( S \) is a \( (m,n) \) tensor, Then we define the tensor product 
    \[ T \otimes S(\omega^{1}, \ldots \omega^{k+m}, V^{1}, \ldots, V^{l+n} ) = \]
    \[ T(\omega^{1},\ldots, \omega^{k}, V^{1}, \ldots, V^{l}) \times S(\omega^{k+1}, \ldots, w^{k+m}, V^{l+1}, \ldots V^{l+n}) \]
    \noindent In general, the tensor product does not commute. 

    \item In index notation, we can write \( T \) as Then in index notation, we can write it as \( T^{\mu_{1}\ldots \mu_{k}}{}_{\nu_{1}\ldots \nu^{l}} \). The basis for the vector space \( T \)
    resides in can be expressed with basis \( \hat{e}_{\mu_{1}} \otimes \ldots \otimes \hat{e}_{\mu_{k}} \otimes \hat{\theta}^{\nu_{1}} \otimes \ldots \otimes \hat{\theta}^{\nu_{l}} \). The components can be found by acting
    the tensor \( T \) on the basis vectors and dual vectors. 

\end{itemize}

\noindent \textbf{Example.} The metric tensor \( \eta_{\mu \nu} \) is a \( (0,2) \) tensor that takes in two vectors and outputs a scalar. The metric acting on two vectors returns the inner product
\[ 
     \eta(V,W) = \eta_{\mu \nu} V^{\mu}W^{\nu} = V \cdot W
\]

\noindent Properties of the inner product should be familiar. Two vectors are orthogonal if their inner product is 0. Note that because the inner product is a scalar, it's invariant under Lorentz transformations. This means that the cartesian
basis is always orthogonal under Lorentz transformation. 
\\

\noindent Additionally, the norm, or "length" of a vector is defined by its inner product with itself \( \eta_{\mu \nu}V^{\mu}V^{\nu} \). Unlike in Euclidean space, there is no restriction for the norm to be strictly positive. A positive
norm is \textbf{spacelike}, a negative norm is \textbf{timelike} and 0 is \textbf{lightlike} or null. 
\\

\noindent The metric \(g_{\mu \nu}\) is an incredibly important tensor. It is symmetric, has nonzero determinant and is used to raise and lower indices. Some of the important roles of the metric include:
    \begin{itemize}
        \item[-] defines path length and proper time
        \item[-] allows us to define a ``shortest distance'' between two points
        \item[-] replaces the Newtonian gravitational field
        \item[-] allows us to define local inertial frames 
    \end{itemize}
    

\noindent We often like to put the metric into a \textbf{canonical form} which means it's a diagonal matrix with entries \(-1,-1 ..., -1, 1..., 1,0 ...0\).
At any point \(p\) on a manifold, there exists a coordinate system \(x^{\hat{\mu}}\) such that \(g_{\hat{\mu}\hat{\nu}}\) is in canonical form and \(\partial_{\hat{\sigma}} g_{\hat{\mu}\hat{\nu}} = 0\). This is called a \textbf{local inertial frame}. This is equivalent to saying
in curved spacetime, we can always find a local inertial frame with the Minkowski metric.
\\


\noindent \textbf{Example.} The Kroneker delta \( \delta^{\mu}{}_\nu \) is a \( (1,1) \) tensor. It inputs a vector and a dual vector and outputs a scalar, but can also be thought of as a map from dual vectors to dual vectors and vectors to vectors.
In this framing, the Kroneker delta is just the identity map. We can define the inverse metric \( \eta^{\mu \nu} \)as satisfying the property
\[ 
    \eta^{\mu \nu} \eta_{{\nu \rho}} = \eta_{\rho \nu}^{\nu \mu} = \delta^{\mu}{}_{\rho}
\]

\noindent Some further properties of tensors:
\begin{itemize}
    
    \item We can contract a vector by summing over a repeated upper and lower index. With this, we can use the metric to raise and lower indices: \( \eta_{\alpha \beta} T^{\alpha} = T_{\beta} \)
    and \( \eta^{\alpha \beta}T_{\alpha} = T^{\beta} \). This is how you would convert a vector to its dual form and vice versa. In euclidean space, the metric is its own inverse so the gradient is actually both a vector and 1-form at the same time.

    \item A \textbf{symmetric} tensor is one where any change of indices leaves the tensor unchanged. For example, if \( S_{\mu \nu \rho} = S_{\nu \mu \rho} \) then we say \( S \) is symmetric in the first two indices. A tensor is \textbf{antisymmetric}
    if it changes sign when indices are changed. \( A_{\mu \nu \rho} = -A_{\nu \mu \rho} \) is antisymmetric in its first two indices.

    \item Symmetric and antisymmetric tensors are very useful in that the contraction of a symmetric and antisymmetric tensor along the corresponding indices is always 0. The symmetric part is \( T_{(\mu\nu)} = \frac{1}{2}(T_{\mu \nu} + T_{\nu \mu}) \) and
    the antisymmetric part is \( T_{[\mu \nu]} = \frac{1}{2}(T_{\mu \nu} - T_{\nu \mu}) \). We sometimes want to exclude indices from being symmetrized or antisymmetrized so we use the vertical bars. For example, to write the antisymmetrized \( R_{\alpha \beta \gamma \delta} \) 
    in only the first and third indices we use \( R_{[\alpha |\beta| \gamma] \delta} = \frac{1}{2}(R_{\alpha \beta \gamma \delta} - R_{\gamma \beta \alpha \delta})\)

    \item For a \( (1,1) \) tensor, we write the trace by leaving off the indices (just the contraction along the repeated index) 
    \[ 
        X = X^{\mu}{}_{\mu} 
    \]
    We can similarly define the trace of a \( (0,2) \) or \( (2,0) \) tensor by raising/lowering into a \( (1,1) \) tensor and then contracting. Note that in this case, the trace is not just the naive sum over the diagonal; the components change when raising/lowering
    the index. The reason behind this is because the sum over the diagonal is not a true scalar; ie it's not invariant under coordinate transform.

\end{itemize}

\section{Riemmannian Geometry}
\subsection{Covariant Derivative}
The partial derivative doesn't behave like a tensor under coordinate transformation, so we need to define a new object, the \textbf{covariant derivative}, that serves as a partial derivative on manifolds and transforms like a tensor. Additionally, it should reduce to the partial derivative in inertial frames. There are
a few motivational details to keep in mind:
\begin{itemize}
    \item The partial derivative \(\partial_\mu\) maps \((k,l)\) tensor fields to \((k, l+1)\) tensor fields. The covariant derivative \(\nabla\) should do the same.

    \item Any good derivative should obey both linearity and the Liebniz rule:
    \[
        \nabla(T + S) = \nabla T + \nabla S
    \]
    \[
        \nabla(T \otimes S) = (\nabla T)\otimes S + T \otimes (\nabla S)
    \]
    \item To obey the Liebniz rule, \(\nabla\) can be written as a partial derivative plus some linear transformation. This gives us the definition of the covariant derivative:
    \[
        \nabla_\mu V^\nu = \partial_\mu V^\nu + \Gamma^\nu_{\mu \lambda}V^\lambda
    \]
    where \(\Gamma^\nu_{\mu \lambda}\) are known as the \textbf{connection coefficients} or \textbf{Christoffel symbols}. They are there to make sure the covariant derivative transforms like a tensor, and they themselves are not tensors because they don't obey the transformation laws.
    
    \item The covariant derivative commutes with contraction, so \(\nabla_\mu(T^\lambda{}_{\lambda\rho}) = (\nabla T)_\mu{}^\lambda{}_{\lambda\rho}\). Additionally, it should reduce to the partial derivative for scalars
    \(\nabla_\mu\phi = \partial_\mu \phi\)

    \item One can show using the last property that the covariant derivative of one forms is given by
    \[
        \nabla_\mu \omega_\nu = \partial_\mu \omega_\nu - \Gamma^\lambda_{\mu \nu}\omega_\lambda
    \]
\end{itemize} 

\noindent Of course, the connection is not unique. Any tensor satisfying the Liebniz rule, linearity, and reducing to the partial derivative can be a covariant derivative, each of which can have a connection associated with it. In general relativity, however, there are several additional constraints:
\begin{itemize}
    \item The connection is \textbf{torsion free}, meaning \(\Gamma^\lambda_{\mu \nu} = \Gamma^\lambda_{\nu \mu}\). This is equivalent to saying the covariant derivative of a scalar field is independent of the order of differentiation, so \(\nabla_\mu\nabla_\nu\phi = \nabla_\nu\nabla_\mu\phi\).

    \item The connection is \textbf{metric compatible}, or \(\nabla_\rho g_{\mu \nu} = 0\). This further means that the covariant derivative commutes with raising and lowering of indices \(g_{\mu \lambda} \nabla_\rho V^\lambda = \nabla_\rho (g_{\mu \lambda}V^\lambda) = \nabla_\rho V_\mu\)

    \item There exists a unique connection satisfying this; the Levi-civita connection
    \[
        \Gamma^\lambda_{\mu \nu} = \frac{1}{2} g^{\lambda \sigma} (\partial_\mu g_{\nu \sigma} + \partial_\nu g_{\mu \sigma} - \partial_\sigma g_{\mu \nu})
    \]
    
\end{itemize}

\subsection{Parallel Transport}
In Euclidean space, we often say a vector doesn't have an associated "position" because we can freely move it around and it will remain unchanged. However, this is not true in curved space. The idea of moving a vector along a path while keeping it constant is called \textbf{parallel transport}. Notably,
parallel transport depends on the path taken between points. 
\\
\\
In fact, there is no ``natural'' way to compare vectors in different tangent spaces, since there are different paths to take. Thus, we can only compare two vectors if they are in the same tangent space. For example, the relative velocities of two particles on different points of a manifold is not well defined!
\begin{itemize}
    \item By "keeping the vector constant," we mean its components must be constant. For an arbitrary tensor, this condition is \(\frac{d}{d\lambda} T^{\mu_1 \mu_2 \cdots \mu_k}{}_{\nu_1 \nu_2 \cdots \nu_l} = \frac{dx^\mu}{d\lambda} \frac{\partial}{\partial x^\mu} T^{\mu_1 \mu_2 \cdots \mu_k}{}_{\nu_1 \nu_2 \cdots \nu_l} = 0\). 
    Then to make this a proper tensor equation, it makes sense to replace the partial derivative with the covariant derivative so we define the \textbf{directional covariant derivative} to be
    \[
        \frac{D}{d\lambda} = \frac{dx^\mu}{d\lambda}\nabla_\mu
    \]
    The equation for parallel transport is then
    \[
        \left(\frac{D}{d\lambda}T\right)^{\mu_1\mu_2\cdots\mu_k}{}_{\nu_1\nu_2\cdots\nu_l} \equiv \frac{dx^\sigma}{d\lambda} \nabla_\sigma T^{\mu_1\mu_2\cdots\mu_k}{}_{\nu_1\nu_2\cdots\nu_l} = 0
    \]
    
    
\end{itemize}

\section{Gravitation}

\subsection{The Newtonian Limit}

\subsection{Stress-Energy Tensor}

We'll be working with fluids in this section. The definition of fluids is not precise
so we'll just say a fluid has friction force negligible in comparison to the pressure.
\begin{itemize}
    \item Fluid particles travelling with four-velocity $\mathbf{u}^{\alpha}$. We can choose 
    to move to inertial frame where $\mathbf{u}^{\mu} = (1,0,0,0)$. This is called the \textbf{comoving frame}.

    \item The \textbf{number current} is $\mathbf{j}^{\mu} = n\mathbf{u}^{\mu} = n(u^0, \mathbf{u}^i)$ where $n$ is the 
    number density (particles per unit volume). The components of the number current is the number of particles
    travelling through unit area per unit time which is the particle flux.

    \item The rest \textbf{mass density} is $\mu_0 = m_0n$ where $m_0$ is the rest mass of a single particle. If there
    is internal energy, then we have $\mu = \mu_0(1+\delta)$ where $\delta$ is the internal energy fraction.
    
    \item The \textbf{momentum density} is $\mathbf{\rho}^i = m_0n\mathbf{u}^{i}$. A nice exercise is to show this is equivalent 
    to the energy flux in special relativity.

\end{itemize}


\noindent We'll now introduce the \textbf{stress-energy tensor} $\mathbf{T}^{ij}$. 

\begin{itemize}
    \item In newtonian mechanics, we have the stress tensor $\mathbf{t}^{ij}$. This represents shear: the ith component of force
    per unit area (pressure) exerted across a surface normal to the jth direction.

    \item For a fluid at rest, the force exerted across a surface is always along the normal and the same in all directions, equivalent
    to pressure $\mathbf{t}^{ij} = p\delta^{ij}$

    \item For a perfect fluid, $\mathbf{T}^{\alpha\beta} = (\rho + p)u^\alpha u^\beta + p\mathbf{g}^{\alpha\beta}$. Careful, $\rho$ is the
    energy density! Meanwhile, $\mathbf{T}^{\alpha}_{\beta} = (\rho + p)u^\alpha u_\beta + p\mathbf{\delta}^\alpha_\beta$

\end{itemize}


\noindent Here, we'll read off some of the physical interpretations of the components of the stress-energy tensor:

\begin{itemize}
    \item $\mathbf{T}^{\alpha}_{\beta} u^\beta = -\rho u^\alpha$. Physically, this quantity's spacial components are the energy flux seen comoving 
    with the fluid. The temporal component $\mathbf{T}^0_\beta u^\beta = -\rho$ is just the rest mass energy density.

    \item $\mathbf{T}^{0}_{i}$ is the flow of the ith component of momentum in the temporal (i = 0) direction. This is the energy flux
    Consequently, $\mathbf{T}^{i}_{0}$ is the energy crossing through the ith surface, so the energy flux. Thus, these two quantities are equivalent.
    In the comoving frame, $\mathbf{T}^{0}_{i} = \mathbf{T}^{i}_{0} = 0$ because the fluid is at rest.

\end{itemize}

\noindent Now, we'll see some results that arise from the Newtonian Limit. Firstly, we assume that the fluid is travelling with low velocity $|u^{\alpha}| << 1$.
Additionally, we have low pressure $\frac{p}{\rho} << 1$ and the speed is non-relativistic $u^{i} \approx v^{i}$.
\begin{itemize}
    \item Firstly, conservation of energy-momentum means that the stress-energy tensor is divergence-free. This gives us
    $\nabla\cdot\mathbf{T} = 0 $ or equivalently $\mathbf{T}^{\alpha\beta}_{;\beta} = 0$. (why???)

    \item Next, reading off the stress-energy tensor components, we have:
        \subitem - $\mathbf{T^{00}} \approx \rho$
        \subitem - $\mathbf{T^{j0}} \approx \rho v^{j}$
        \subitem - $\mathbf{T^{jk}} \approx \rho v^j v^k + p \delta^{jk}$
    
    \item The continuity equation $\frac{\partial{\rho}}{\partial{t}} + (\rho v^k)_{;k} = 0$ effectively states conservation of mass. The rate of change of density 
    plus the net mass flow out is 0. Combining with the divergence free condition, we obtain
    $\frac{\del v^j}{\del t} + (v \cdot \nabla)v = -\frac{1}{\rho}\nabla p$. This is the \textbf{Euler Equation}.

    \item The equations of state relate the pressure and the energy density for a fluid. There are several possible cases:
        \subitem - Vacuum : $p = \rho = 0$ (Inflationary universe)
        \subitem - Dust : $p = 0$  (FLRW universe)
        \subitem - Fluid : $p = \gamma \rho$
        \subitem - Polytropic : $p = n\rho ^{(1 + \frac{1}{n})}$ where n is the polytropic index.

\end{itemize}


\subsection{Gravitational Waves}
\noindent Recall the form of the metric in the linearized theory $g_{\mu\nu} = \eta_{\mu\nu} + h_{\mu\nu}$.
\begin{itemize}
    \item The einstein equations give us $2(R_{\alpha\beta} - \frac{1}{2}g_{\alpha\beta}R) = 16\pi G T_{\alpha\beta}$
    \item Introducing the trace-reversed metric perturbation, $\bar{h}_{\alpha\beta} \equiv h_{\alpha\beta} - \frac{1}{2}h$
    and expanding the einstein equation, we can simplify this to $\Box \bar{h}_{\alpha\beta} = \bar{h}_{\alpha\beta ,\mu}^{,\mu} = -16\pi G T_{\alpha\beta}$.
    The box is the d'Alembert operator and is the equivalent of the laplacian in spacetime. This nicely gives us a wave equation for $\bar{h}$.

    \item The solution for this wave equation is $\bar{h}_{\mu\nu} = Re(A_{\mu\nu}e^{i k_{\alpha}x^{\alpha}})$ where $A_{\mu\nu}$ gives the amplitude and $k_{\alpha}$
    is the wavevector $(w,\mathbf{k})$. From the d'Alembert operator, we have \( k_{\alpha}k^{\alpha} = 0\), or \( k \) is a null vector. This means gravitational waves
    propogate at the speed of light.

    \item We can choose a Lorentz Gauge to simplify this further (what does this mean?) with \( \overline{h_{\mu\nu}^{,\nu}} = 0 \). This gives us \( A_{\mu\alpha}k^{\alpha} = 0 \)
    meaning these are transverse waves (why?)

    \item Starting with the Symmetric 4x4 tensor, we have 10 independent coordinates. Being tranverse and traceless takes away 4 degrees of freedom (why?) and the gauge condition
    takes away 4 more so we are left with 2 polarizations. Imposing all conditions gives us the matrix 
    \[ 
        \begin{pmatrix}
        0 & 0 & 0 & 0 \\
        0 & a & b & 0 \\
        0 & b & -a & 0 \\
        0 & 0 & 0 & 0 \\
        \end{pmatrix}
    \]
    
    \item Both linear and circular polarization like light (This section needs work)

    
\end{itemize}

\noindent Now we'll derive the metric from a weakly gravitating source. Suppose we have a static point mass M located at point \( r = 0 \) at all times in the weak field limit.
In the body's frame we have \( |h_{\mu\nu}| \ll 1 \). 
\begin{itemize}
    \item Because the body is stationary, the stress energy tensor has components \( T_{00} = \delta(x)M \), and all other 
    components are 0. 
    \item From the einstein equation and poisson equations, we have \(  \overline{h}_{00} = \frac{4M}{r}\) and all other components are 0. 
    \item Converting from \( \overline{h} \) back to \( h \) and substituting into the equation for the metric \( ds^{2} = (\eta_{\mu\nu} + h_{\mu\nu})dx^{\mu}dx^{\nu} \) we 
    obtain the metric:
    
    \[ 
        ds^{2} = 1 \left( 1 - \frac{2M}{r}\right)dt^{2} +\left(1+\frac{2M}{r}\right)(dx^{2} + dy^{2} + dz^{2})
    \]

    

\end{itemize}





\section{Black Holes}
Rev. John Mitchell found in 1783 that an object with $500_{\odot}$ but of same density will have its escape velocity to be greater than the speed of light. Then 
"all light emitted from such a body would be made to return towards it."
\begin{itemize}
    \item Every mass has a radius such that if the mass is small enough than this radius, it will collapse into a black hole. This is the \textbf{Schwarzchild Radius}. 
    $R_{s} = \frac{2 G M}{c^2}$

\end{itemize}


\subsection{Schwarzchild Metric}
The Schwarzchild Metric has two forms:
\begin{itemize}
    \item The timelike form $d\tau^{2} = (1 - \frac{2M}{r})dt^2 - \frac{dr^2}{(1 - \frac{2M}{r})} - r^2(d\theta^2 + \sin^2\theta d\phi^2)$
    \item The spacelike form $d\tau^{2} = (1 - \frac{2M}{r})dt^2 - \frac{dr^2}{(1 - \frac{2M}{r})} - r^2d\Omega^2$
    \item We can see the metric reduces to flat spacetime in the case where $r \rightarrow \infty$ and $m \rightarrow \infty$
    \item Upon examining the temporal component of the metric, we that its less than that of the temporal component in flat spacetime, so this properly predicts gravitational
    redshift. Similarly, the spacial component is larger than that in flat spacetime, so there's spacetime curvature.
\end{itemize}


\subsection{Coordinates Systems}
These are three commonly used frames:
\begin{itemize}
    \item Free-Float Frame: Frame free-floating in space, think unpowered spacecraft. Locally its flat minkowski space. For observers on shells, they see spacecraft infalling towards
    the black hole.
    \item Spherical-Shell Frame: Frame standing still on a shell of fixed radius (eg: Earth's frame). For small time lapses and regions, special relativity applies. Notice that inside a
    black hole, $dr$ becomes time like while $dt$ becomes spacelike. This means staying on a shell at constant radius is impossible, so this frame cannot be used to describe events inside 
    a black hole.
    \item Schwarzchild Frame: Global book-keeping coordinates. Idea is that neither spherical shell frame nor free float frame can know all information so both send information to a Global
    frame that puts it together.
\end{itemize}

\noindent A common result from black holes is that from the perspective of someone outside of the black hole, we see an object travel into the black hole past the event horizon. But from the perspective
of the falling object, it takes infinite time to reach the event horizon!
\begin{itemize}
    \item We can calculate the Riemman tensor in the shell frame: $R_{trtr} = -\frac{2M}{r^3} = -R_{\theta\phi\theta\phi}$ and 
    $R_{t\theta t\theta} = R_{t \phi t \phi} = \frac{M}{r^3} = -R_{r\theta r \theta} = -R_{r \phi r\phi}$
    \item Meanwhile in the free falling frame we have $R_{trtr} = -\frac{2M}{r^3}(1-\frac{2M}{r})^{-1}$. This isn't singular at $r = 2M$ but it is at $r=0$. The curvature invariant (also known as
    the Kretschmann scalar) $R = R_{trtr}R^{trtr} = \frac{48M^2}{r^6}$. This gives a measure of the strength of the gravitational field using curvature and this also blows up to infinity at $r = 0$.
\end{itemize}

\subsection{Orbits}
Because the Schwarzchild metric is time independent and spherically symmetric, we must have killing vectors $\xi_{t} = (1,0,0,0)$ and $\xi_{\phi} = (0,0,0,1)$. 
For any four-velocity $u$, $\xi \cdot u$ is constant. combining this with $p^{\mu} =m u^{\mu}$, we obtain $\frac{E}{m_0 c^2}$ is constant along a path.

\begin{itemize}
    \item Radial Infall: Let $d\tau$ be the proper time for an observer at distance r from the center of a black hole (shell frame). 
    At $r \rightarrow \infty$, we have $\frac{E_\infty}{m_0 c^2} = \frac{dt}{d\tau}$. But $E_\infty$ is a constant. So any trajectory starting at rest from $r = \infty$ will have
    same energy per unit mass throughout. This gives us $d\tau ^2  = (1- \frac{2M}{r})^2 dt^2$.

    \item With this expression we can obtain the velocity recorded by a Schwarzchild observer:   
    \newline$\frac{dr}{dt} = -(1-\frac{2M}{r})(\frac{2M}{r})^{\frac{1}{2}}$. At $r = 2M$, this goes to 0 so we see that the bookeeper observes the particle slowing down as it approaches the event horizon,
    taking infinite time to reach there.

    \item If we instead moved to the frame of the shell observer, we can obtain from the Schwarzchild metric: $\frac{dr_{shell}}{dt_{shell}} = -(\frac{2M}{r})^\frac{1}{2}$. For an observer at the event horizon,
    they would see the object falling at the speed of light! The energy observed also depends on the distance from the center (while $E_{\infty}$ is constant, $E_{shell}$ is not).

    \item Meanwhile, for the free falling frame, we can actually compute the time it takes for the observer to reach $r = 0$. Using the chain rule, $\frac{dr}{d\tau} = \frac{dr}{dt}\frac{dt}{d\tau}$ we can obtain 
    $\tau = \frac{4}{3}M$. In seconds this is $\tau_{second} = \frac{4}{3}\frac{M}{M_{odot}}(5\mu sec)$. Interestingly, the time scales linearly with mass, so for larger black holes it actually takes more time to
    fall! For reference, Sagittarius A, the supermassive black hole in the center of the galaxy, it takes roughly 1 minute to fall. Meanwhile, TON618, one of the largest black holes takes 11 days to fall to the center.

\end{itemize}

\noindent Now we'll go over orbits. In the case with $r > 20R_{s}$, the newtonian approximation is still accurate. For $r \leq 6M$, no free orbits are possible. The mass spirals in but still may exert enough energy to escape.
Beyond $r \leq 2M$, no escape is possible. Inside the horizon, as mentioned before, space and time components are interchanged. Travelling to smaller r is as "inevitable" as outside a black hole we travel to greater t.

\begin{itemize}
    \item Recall the conserved quantity from earlier $\xi_t \cdot u^t$. We call this the "energy per unit mass" and we'll just set this equal to $\frac{E}{m}$. The free fall case assumes this is one but this is not true generally.
    Write the angular momentum per unit mass as $\frac{L}{m} = r^2 \frac{d\phi}{d\tau}$. Using these, we have the equation of motion 
    \[ 
        \left(\frac{dr}{d\tau}\right)^2 = \left(\frac{E}{m}\right)^2 - \left(1 - \frac{2M}{r}\right)\left[1 + \frac{(L/M)^2}{r^2}\right] \equiv \left(\frac{E}{m}\right)^2 - \left(\frac{V}{m}\right)^2 
    \].
     This gives us an effective potential of 
    \[
        \left(\frac{V(r)}{m}\right)^2 \equiv \left(1 - \frac{2M}{r}\right)\left[1 + \frac{(L/m)^2}{r^2}\right]
    \]
    In the Newtonian case, we have an elliptical orbit. But for black holes, orbits tend to spend more time at smaller radii, the additional time near the inner edge of the orbit is the "\textbf{dwell time}."
    
    \item Stable circular orbits at $V_{min}$ and unstable orbits occur at $V_{max}$. Beyond $V_{max}$ the orbit precesses because of the dwell time. (this section needs work)
\end{itemize}

\subsection{Massless Particle}
We'll now cover orbits of massless particles. This is especially important for describing the behavior of light near a black hole. The procedure for massive particles won't work here for several reasons.
    \begin{itemize}
        \item The spacetime interval for massless particles is 0 (why?). Solving for the coordinate velocity of light in these coordinates gives us $r\frac{d\theta}{dt} = \pm (1-\frac{2M}{r})^\frac{1}{2}$ tangentially and
        $\frac{dr}{dt} = \pm (1 - \frac{2M}{r})$ radially. This looks like it
        means the speed of light can even approach 0 but this actually isn't a problem because this interpretation is unphysical. Even if the Schwarzchild observer "sees" the speed of light as this, any shell observer still 
        sees it at constant c.
        \item At $r = \infty$, the special relativity case we have $p = p_{\infty}$. The angular momentum L can be written using this momentum as $L = bp_{\infty}$ where b is known as the \textbf{impact parameter}. (Why?)
        Using the expression for $P$ in special relativity, we can write $b = \frac{L/m}{E/m}$

        \item The impact parameter conveniently allows us to evaluate the speed even with 0 mass: 
        \newline The radial velocity: $(\frac{dr}{dt})^2 = (1 - \frac{2M}{r})^2 - (1 - \frac{2M}{r})^3\frac{b_{photon}^2}{r^2}$ 
        \newline The tangential velocity : $r\frac{d\phi}{dt} = \frac{b_{photon}}{r}(1 - \frac{2M}{r})$
        \newline Thus the total speed is $v_{photon} = (1 - \frac{2M}{r})[1 + \frac{2Mb^2}{r^3}]^\frac{1}{2}$.
        \newline With a similar technique, we can verify that in the shell coordinates, the speed of the particle is constant 1.

        \item The equivalent equation of motion in the shell frame is $\frac{1}{b^2}(\frac{dr_{shell}}{dt_{shell}})^2 = \frac{1}{b^2} - \frac{(1 - \frac{2M}r)}{r^2}$. $\frac{1}{b^2}$ plays the role of the energy squared while
        the second term is the equivalent of the effective potential.

        \item This potential has no stable orbits. A photon may enter a circular "knife edge" orbit before either escaping or plunging.
    \end{itemize}

\section{Cosmology}
\subsection{The Present Universe}
In the last 100 years, there were many important observations about our current universe.
\begin{itemize}
    \item Hubble discovered the universe is expanding in 1929. The speed at which galaxies recessed from us is proportional to their distance $v = Hr$ where H is the hubble constant.
    \item In 1965, Penzias and Wilson discovered the cosmic microwave background has temperature $T = 273^{\circ} K$. It is also highly isotropic and homogeneous. 
    \item Matter density upper bounds at $10^{-29} g/cm^3$ and lower bound at $5 \times 10^{-31}g/cm^3$. The Baryon number Photon number ratio is $10^-{8} - 10^{-10}$
    \item From stars, galaxies and globular clusters, the Helium-4 mass fraction is 0.2-0.3, while Deutereum is $3 \times 10^{-5}$
    \item There is a high degree of local structure, with galaxies, galaxy clusters and galaxy filaments. Supermassive objects like galaxy filaments, voids and walls
    put the cosmological principle into question.
    \item Dark energy (or vacuum energy) is at 0.69, corresponding to $69\%$ of total energy density in the universe.
\end{itemize}

\noindent Now we'll derive the Einstein Equations for a Robertson-Walker Universe, and show that this satisfies the isotropy and homogeneity conditions at all points.
\begin{itemize}
    \item Consider an infinite, homogeneous gas cloud with mass density $\rho$. Choose an arbitrary point $A$ as the origin and consider a thin spherical shell centered around $A$ with 
    radius $r$. By conservation of energy, we have $(\frac{\dot{r}}{r})^2 + \kappa(r) = \frac{8G\rho}{3}$. $\kappa$ is $\frac{k}{r^2}$ where k is the total energy.
    \item By using the First Law of Thermodynamics and assuming constant entropy, we can obtain $\frac{d\rho}{dt} + \frac{3}{r}\frac{dr}{dt}(\rho + p) = 0$. Taking the first derivative of
    the previous line and combining, we have $\frac{\ddot{r}}{r} + \frac{4\pi G}{3}(\rho + 3p) = 0$. 
    \item Choosing a coordinate comoving with the gas cloud, we can define $\mathbf{r(t)} = a(t)\mathbf{r_0}$. $a(t)$ is the scale factor that depends only on time (no dependence on space because isotropy).
    From this, we can rewrite our earlier equations to obtain
    \newline $\frac{\ddot{a}}{a} + \frac{4\pi G}{3}(\rho + 3p) = 0$
    \newline $(\frac{\dot{a}}{a})^2 + \frac{\kappa(r)}{a^2} = \frac{8G\rho}{3}$. These are the the einstein equations in component form. 
    \item The continuity equation of the stress energy tensor can similarly be rewritten $\dot{\rho}+3(\frac{\dot{a}}{a})(\rho + p) = 0$.
    \item These equations can be further simplified using equations of state. 
    \newline Dust: $p = 0$
    \newline Radiation: $p = \frac{1}{3}\rho$
    \newline Vacuum Energy: $p = -\rho$
\end{itemize}

\subsection{History of the Universe}
\noindent Hubble's discovery of the expansion of the universe and Penzias \( \& \) Wilson's discovery of the cosmic microwave background suggest the existence of the universe in a "hot, dense" state in
the very beginning, before expansion. This is the popularly named "\textbf{Big Bang}" model (Gamow 1948, Dicke, Peebles 1965).
\\
\\
\noindent Furthermore the existence of General Relativity combined with the energy dominance condition on matter (energy is positive, and energy-momentum doesn't travel faster than light) implies the existence
of a cosmological singularity (Penrose, Hawking, Geroch 1966). (Why?)

\begin{itemize}
    \item Starting with the Friedmann universe, we can do a convenient change of variables by replacing the cumbersome \( (1-kr^{2})^{{-1}} \) with \( \chi \). This is known as the comoving distance. As the universe
    expands, the proper distance between locations change with time but \( \chi \) remains constant \( D(t) = a(t)\chi \).
    
    \[ 
        ds^2 = -dt^{2} + a^{2}(t)[d\chi^{2} + f^{2}(\chi)(d\theta^{2} + \sin^{2}{\theta}d\phi^{2})] 
    \]

    where \( f(\chi) = r \). This gives us nice forms for \( \chi \):
    \begin{enumerate}[label = \roman*]
        \item \( k = 1 \) Closed universe: \( f(\chi) = \sin{\left( \chi \right)} \)
        \item \( k = 0 \) Flat universe: \( f(\chi) = \sinh{\chi} \)
        \item \( k = -1 \) Open universe: \( f(\chi) = \chi \)
    \end{enumerate}
    
    
    In the case where \(f(\chi) = \sin{\left( \chi \right)} \) and \( 0 \leq \chi \leq 2\pi \) we have a closed universe. For \( f(\chi) = \chi \) and \( 0 \leq \chi \le \infty \) this
    is a flat universe. Finally, for \( f(\chi) = sinh(\chi) \) then this is an open universe. (What do these mean?)
    

    \item Substituting the equatins of state into the continuity equation we can obtain scaling factors of the energy density with \( a \)
    In the flat universe,

    \begin{enumerate}[label = \roman*]
        \item dust : \( \rho_{m} \sim a^{-3} \)
        \item radiation : \( \rho_{\gamma} \sim a^{-4} \)
        \item vacuum : \( \rho \sim const \)
    \end{enumerate}

    \item Solving the einstein equation with the Friedmann metric we obtain the following:
    
    \[ 
        G_{00}  \left(\frac{\dot{a}}{a}\right)^{2} + \frac{k}{a^{2}} - \frac{\Lambda}{3} = \frac{8\pi G \rho}{3}
    \]
    This is the energy conservation equation, which relates the rate to energy density and

    \[ 
        G_{{ii}} \to 3\left(\frac{\ddot{a}}{a}\right) = -4 \pi G(\rho + 3p) 
    \]
    This is the acceleration equation. 

    \item If we solve these equations using the equations of state in the flat case (\( k = 0 \)) we obtain the following:
    \begin{enumerate}[label = \roman*]
        \item In a matter dominated universe: \( a \sim t^{\frac{2}{3}} \)
        \item Radiation dominated: \( a \sim t^{\frac{1}{2}} \)
        \item Vacuum dominated: \( a \sim e^{Ht} \)
    \end{enumerate}

\end{itemize}

\subsection{Singularity}
A big bang corresponds to a scale factor of \( a = 0 \). We can ask the question if our universe was ever at this point ie: 
does it take finite time to reach \( a = 0 \) in the past?
\begin{itemize}
    \item First, assume the universe is radiation-dominated (which it would've been at that time) and with \( \Lambda = 0 \)
    Then from the \( G^{0}_{0} \) equation, we obtain \( a^{2} = \left(\frac{32}{3}\pi G B\right)^{\frac{1}{2}}(t - t_{0}) \)
    where \( \rho = Ba^{-4} \). This means it takes finite time to reach \( a=0 \)
    \item A similar analysis for the \( G^{i}_{i} \) equation gives us \( \frac{\dot{a}}{a} \to \infty \) at \( t = t_{0} \),
    which implies this time has infinite expansion rate.
    \item Causality laws: the speed of sound is always less than the speed of light imply that \( \rho \to \infty \) at this time.
\end{itemize}
Since the universe began at \( R = 0 \) a finite time in the past, this makes the Big Bang inevitable. This point in time is called
the cosmological singularity. Penrose and Hawking's singularity theorems (1973) show that the universe must have been a singularity
even if the conditions of isotropy and homogeneity were relaxed.


\subsection{The FLRW Universe}

\begin{itemize}
    \item We define the proper distance between 2 galaxies as the integral
    \[ 
        d(t) = a(t) \int^{r}_{0}\frac{dr}{\sqrt{1-kr^{2}}} = a(t) 
        \begin{cases}
        \sin^{-1}{r} & \text{for } k = 1 \\
        r & k = 0 \\
        \sinh^{-1}{r} & k = -1
        \end{cases}
    \]

    \noindent Meanwhile the proper velocity is 
    \[ 
        v(t) \equiv \dot{d}(t) = \frac{\dot{a}(t)}{a(t)}d(t) = H(t)d(t) 
    \]

    \noindent this is Hubble's law, where \( H_{0}(t_{0}) \) is the present value of the hubble constant.


    \item Let \( r_{ph} \) be the coordinate of the most distant object observable, with the earliest light coming from time \( t_{min} \).
    For light rays, we can set \( ds = 0 \) giving us 
    \[ 
        \chi_{p} = c\int^{t_{0}}_{t_{min}} \frac{dt}{a(t)} = \int^{r_{ph}}_0\frac{dr}{\sqrt(1-kr^{2})}
    \]
    The closed universe model has \( r_{ph} = \sin(\chi_{p}) \) which is a finite value, so there exists a horizon beyond which you cannot see
    past (why?). This is called the particle horizon.

\end{itemize}

\subsection{The CMB}

\subsection{The Early Universe}









\end{document}